\section{Physical interpretation: window bounds and classical/quantum non-overlap}
\label{sec:physical_interpretation}

\subsection{Why does the quantum Efimov window exist for
  \texorpdfstring{$\lambda_{\rm sp}\in(5.1,\,13.2)\,l_{\rm Pl}$}{lambda in (5.1,13.2) l\_Pl}?}
\label{sec:why_window}

The quantum Efimov window $C_{\rm Pl}\in\bigl(C_Q^{(3B)},\,C_Q^{(2B)}\bigr)$
exists for scalar range $\lambda_{\rm sp}\in(5.1,\,13.2)\,l_{\rm Pl}$.
Each boundary has a clear physical origin.

\paragraph{Upper boundary $\lambda_{\rm sp}^{\rm upper}=13.2\,l_{\rm Pl}$
  ($m_{\rm sp}=0.076\,l_{\rm Pl}^{-1}$):}
As the scalar range $\lambda_{\rm sp}$ increases, the Yukawa two-body
potential becomes more long-range and the two-body quantum threshold
$C_Q^{(2B)}(\lambda_{\rm sp})$ decreases.  At
$\lambda_{\rm sp}=13.2\,l_{\rm Pl}$ the threshold reaches the Planck
coupling:
\begin{equation}
  C_Q^{(2B)}(m_{\rm sp}=0.076) = C_{\rm Pl} = 0.282.
\end{equation}
For $\lambda_{\rm sp}>13.2\,l_{\rm Pl}$ (longer range), pairs
quantum-mechanically bind at $C_{\rm Pl}$, so the three-body cluster
is no longer a ``Efimov-type'' purely three-body state —
pairs already bind by themselves.

\textbf{Physical meaning:} $\lambda^{\rm upper}$ is the Yukawa range
at which the Planck coupling first becomes strong enough to form
quantum-mechanically bound \emph{pairs}.  Above this range, two-body
physics dominates.

\paragraph{Lower boundary $\lambda_{\rm sp}^{\rm lower}=5.1\,l_{\rm Pl}$
  ($m_{\rm sp}=0.198\,l_{\rm Pl}^{-1}$):}
As the scalar range decreases (heavier mediator), the three-body
Feynman integral $I_Y(d,d,d;\,m_{\rm sp})$ becomes more rapidly
screened for $d\gg\lambda_{\rm sp}$, reducing the range over which
three-body attraction acts.  The three-body quantum threshold
$C_Q^{(3B)}$ therefore increases.  At $\lambda_{\rm sp}=5.1\,l_{\rm Pl}$:
\begin{equation}
  C_Q^{(3B)}(m_{\rm sp}=0.198) = C_{\rm Pl} = 0.282.
\end{equation}
For $\lambda_{\rm sp}<5.1\,l_{\rm Pl}$ (shorter range), the Planck
coupling is insufficient to bind a quantum three-body cluster —
the zero-point kinetic energy $3/(8d^2)$ in the collective coordinate
overcomes the screened attraction.

\textbf{Physical meaning:} $\lambda^{\rm lower}$ is the Yukawa range at which
the three-body TGP interaction becomes too short-ranged to support a
quantum bound cluster at $C_{\rm Pl}$.  Below this range, zero-point
pressure wins.

\paragraph{Summary of window origin.}
The window $\lambda_{\rm sp}\in(5.1,\,13.2)\,l_{\rm Pl}$ is the range
in which:
\begin{itemize}
\item The scalar is \emph{long-range enough} that three-body attraction
  overcomes zero-point pressure (requires $\lambda_{\rm sp}>5.1\,l_{\rm Pl}$).
\item The scalar is \emph{short-range enough} that pairs remain
  quantum-mechanically unbound at $C_{\rm Pl}$
  (requires $\lambda_{\rm sp}<13.2\,l_{\rm Pl}$).
\end{itemize}
The window width $\Delta\lambda = 8.1\,l_{\rm Pl}$ (factor $\times2.6$ in
$m_{\rm sp}$) is not fine-tuned.

\paragraph{Connection to wavefunction (ex29).}
The ground-state wavefunction at $C_{\rm Pl}$, $m_{\rm sp}=0.1$ peaks at
$d_{\rm peak}=3.5\,l_{\rm Pl}$ with $\langle d\rangle=4.5\,l_{\rm Pl}$.
This is well within the Yukawa range $\lambda=10\,l_{\rm Pl}$, explaining
why the three-body force is effective: the cluster lives where
$I_Y(d,d,d;\,m_{\rm sp})$ is large.  At the lower boundary
($\lambda=5.1\,l_{\rm Pl}$, $m_{\rm sp}=0.198$), the wavefunction would need
to sit at $d\lesssim\lambda=5.1\,l_{\rm Pl}$ to feel the three-body force,
but zero-point energy pushes it outward — the potential well is too narrow
to contain the state (from section~C of ex29: $\langle d\rangle$ grows
from $4.2$ to $8.1\,l_{\rm Pl}$ as $m_{\rm sp}$ approaches $0.198$).

\subsection{Why are the classical and quantum windows non-overlapping?}
\label{sec:nonoverlap}

The central result of Sec.~\ref{sec:quantum_window_msp} is that the
classical Efimov window for $C_{\rm Pl}$ (at $m_{\rm sp}\in(0.235,\,0.38)$)
and the quantum Efimov window (at $m_{\rm sp}\in(0.076,\,0.198)$) are
\emph{mutually exclusive}.  We now explain why.

\subsubsection{The role of zero-point energy}

The classical and quantum treatments of the collective coordinate $d$ differ
in how zero-point kinetic energy enters:

\begin{itemize}
\item \textbf{Classical:} The system minimises $E_{\rm eff}(d)=\frac{3}{8d^2}
  +V_2+V_3$ over $d$.  The ZP term acts as a static repulsive ``pressure''
  at small $d$.  Binding requires $\min_d E_{\rm eff}<0$.

\item \textbf{Quantum:} The operator $\hat{T}=-\frac{1}{2\mu}\frac{d^2}{dd^2}$
  contributes kinetic energy that depends on the wavefunction curvature.
  The ground-state energy $E_0 = \langle\psi|\hat{H}|\psi\rangle$ includes
  both potential \emph{and} kinetic contributions.  Since $\psi$ is spread
  over $\sigma_d\sim 2\,l_{\rm Pl}$ (ex29), it samples a range of $V_{\rm eff}$
  values around the minimum, raising the effective energy.
\end{itemize}

The net effect is that the quantum criterion for binding ($E_0<0$) requires
a \emph{deeper and broader} potential well than the classical criterion
($E_{\rm eff}^{\rm min}<0$).  This raises both thresholds:
$C_Q^{(3B)}>C_{\rm cl}^{(3B)}$ and $C_Q^{(2B)}>C_{\rm cl}^{(2B)}$,
and the quantum window is shifted upward in $C$ by $\approx0.063$ at
$m_{\rm sp}=0.1$.

\subsubsection{Why the shift causes non-overlap}

At fixed $C=C_{\rm Pl}=0.282$:

\begin{enumerate}
\item \textbf{Classical window} ($m_{\rm sp}\in(0.235,\,0.38)$, short-range
  Yukawa, $\lambda\in(2.6,\,4.3)\,l_{\rm Pl}$):
  Here $C_{\rm cl}^{(3B)}<C_{\rm Pl}<C_{\rm cl}^{(2B)}$.
  The classical equilibrium separation is $d_{\rm cl}\sim1.5$--$2\,l_{\rm Pl}$,
  well within the Yukawa range $\lambda$.
  However, at these short Yukawa ranges the potential well is very narrow
  (width $\Delta d\sim2\,l_{\rm Pl}$) and shallow.
  The quantum ground-state wavefunction in this shallow well has
  $\sigma_d\sim\Delta d$, so kinetic energy is large and
  $E_0 = E_{\rm eff}^{\rm min} + \langle\hat{T}\rangle \gg 0$ —
  zero-point energy destroys the binding.
  In ex29 at $m_{\rm sp}=0.18$ (close to the classical window):
  $\langle d\rangle=8.1\,l_{\rm Pl}\gg\lambda=5.6\,l_{\rm Pl}$ and
  the three-body force is effectively absent.

\item \textbf{Quantum window} ($m_{\rm sp}\in(0.076,\,0.198)$, longer-range,
  $\lambda\in(5.1,\,13.2)\,l_{\rm Pl}$):
  Here $C_Q^{(3B)}<C_{\rm Pl}<C_Q^{(2B)}$.
  The potential well is broader (extends over $d\sim1$--$6\,l_{\rm Pl}$
  within the Yukawa range), so the wavefunction can be contained while
  zero-point energy remains manageable.
  The quantum bound state exists precisely because the broader range allows
  $\langle d\rangle\sim4.5\,l_{\rm Pl}\lesssim\lambda=10\,l_{\rm Pl}$
  (ex29, $m_{\rm sp}=0.1$).
  However, classically the pair threshold $C_{\rm cl}^{(2B)}\sim0.18$
  is already below $C_{\rm Pl}$, so pairs \emph{classically} bind —
  the classical Efimov-type 3B-only window at $C_{\rm Pl}$ does not exist
  in this range of $m_{\rm sp}$.
\end{enumerate}

\subsubsection{Schematic}

\begin{table}[h]
\centering
\caption{Physical mechanism behind classical/quantum window separation.}
\label{tab:nonoverlap_mechanism}
\begin{tabular}{p{3cm}p{4cm}p{4.5cm}p{2.5cm}}
\hline
$m_{\rm sp}$ range & $\lambda_{\rm sp}$ & Classical binding & Quantum binding \\
\hline
$<0.076$ ($\lambda>13.2$) & Long-range & pairs + triples bind & pairs bind at $C_{\rm Pl}$ \\
$0.076$--$0.198$ ($\lambda=5$--$13$) & Medium & pairs bind classically, no 3B-only window & 3B-only window: \textbf{quantum Efimov} \\
$0.198$--$0.235$ ($\lambda=4.3$--$5.1$) & Medium-short & pairs bind classically & unbound \\
$0.235$--$0.38$ ($\lambda=2.6$--$4.3$) & Short & 3B-only classical Efimov & unbound: ZP destroys it \\
$>0.38$ ($\lambda<2.6$) & Very short & neither binds & neither \\
\hline
\end{tabular}
\end{table}

\subsubsection{Key insight: WKB invalidity as a diagnostic}

A quantitative signature of the non-overlap is the WKB ratio
$\lambda_{\rm dB}/L_{\rm well}$ at the quantum ground state (ex29):
\begin{equation}
  \frac{\lambda_{\rm dB}}{L_{\rm well}} = \frac{16.0}{5.3} \approx 3.0
  \quad (m_{\rm sp}=0.1,\; C=C_{\rm Pl}),
\end{equation}
which is $\gg1$, confirming deep quantum mechanics.  In the classical Efimov
regime (short-range Yukawa), the well would be narrow enough that the
WKB ratio approaches unity — but the state is then pushed out by zero-point
energy.  The two regimes are incompatible at $C_{\rm Pl}$.

\subsection{Quantum state characterisation at
  \texorpdfstring{$C=C_{\rm Pl}$}{C=C\_Pl}}
\label{sec:wavefunction_summary}

The ground state at $C_{\rm Pl}=0.282$, $m_{\rm sp}=0.1\,l_{\rm Pl}^{-1}$
(ex29) has the following properties:

\begin{table}[h]
\centering
\caption{Ground-state wavefunction properties at $C_{\rm Pl}$, $m_{\rm sp}=0.1$.}
\label{tab:psi0_summary}
\begin{tabular}{ll}
\hline
Quantity & Value \\
\hline
$E_0$ & $-0.0696\,E_{\rm Pl} = -1.36\times10^8\,\text{J}$ \\
$d_{\rm peak}$ & $3.51\,l_{\rm Pl} = 5.45\times10^{-35}\,\text{m}$ \\
$\langle d\rangle$ & $4.52\,l_{\rm Pl}$ \\
$\sigma_d$ & $2.10\,l_{\rm Pl}$ \\
$\sigma_d/\langle d\rangle$ & $0.46$ (broad, non-classical) \\
$d_{\rm cl}$ (classical minimum) & $1.37\,l_{\rm Pl}$ \\
$\langle d\rangle/d_{\rm cl}$ & $3.30$ (quantum delocalisation) \\
$\lambda_{\rm dB}/L_{\rm well}$ & $3.03$ (WKB invalid) \\
Tunnelling fraction & $20.8\%$ \\
$V_3/V_2$ at $\langle d\rangle$ & $2.74$ (3-body dominant) \\
Number of bound states & 2 \\
\hline
\end{tabular}
\end{table}

\noindent
The second bound state ($n=1$) has $E_1=-6.5\times10^{-4}\,E_{\rm Pl}$
and peaks at $d_{\rm peak}^{(1)}=13.3\,l_{\rm Pl}$ — a very diffuse,
weakly-bound excited state analogous to the first Efimov excited state
in nuclear/atomic physics.  The ratio of spatial scales
$d_{\rm peak}^{(1)}/d_{\rm peak}^{(0)}\approx3.8$ is much smaller than
the classical Efimov ratio $e^{\pi/s_0}\approx22.7$, confirming that TGP
does not produce an infinite geometric tower.  The Yukawa screening
limits the well depth, allowing at most two bound levels at $C_{\rm Pl}$.

As $m_{\rm sp}$ approaches the upper boundary $0.198\,l_{\rm Pl}^{-1}$,
the bound state becomes increasingly diffuse
($\langle d\rangle$ grows from $4.2$ to $8.1\,l_{\rm Pl}$,
tunnelling from $20\%$ to $32\%$) — a classic signature of a
state approaching the dissociation threshold.

\subsection{Summary}

\begin{enumerate}
\item The quantum Efimov window bounds at $\lambda\approx5$ and
  $13\,l_{\rm Pl}$ are set by the two-body and three-body quantum
  thresholds coinciding with $C_{\rm Pl}$ — not by any fine-tuning.
\item The classical and quantum windows are non-overlapping because the
  zero-point energy operator raises both thresholds by $\approx0.063$ in $C$,
  shifting the quantum window to longer Yukawa ranges where the potential
  well is broad enough to contain the wavefunction.
\item The quantum ground state at $C_{\rm Pl}$ is deeply non-classical:
  $\langle d\rangle/d_{\rm cl}=3.3$, WKB ratio $=3.0$, tunnelling $=21\%$.
  Classical methods give qualitatively wrong predictions for the existence
  and location of bound states.
\item A second bound state exists at $m_{\rm sp}=0.1$, $C=C_{\rm Pl}$,
  at $d\sim13\,l_{\rm Pl}$.  This analogue of the first Efimov excitation
  has no infinite tower — Yukawa screening limits the count to $\leq2$.
\end{enumerate}
